% DRAFT (LLM-drafted from claims ledger rows J2, F5, F10, F14, M1-M5, P6, P7; the author reviews and confirms).
\section{Methodology findings}
\label{sec:findings}

The protocol was designed to avoid known pitfalls; applying it surfaced others. Each finding below is measured on
our data, and several are mistakes we made first and caught later.

\subsection{Probability rounding makes log loss unusable against one API}
\label{sec:find-rounding}
\jev{} returns probabilities rounded to two decimals, so the gold option often receives exactly zero when the model
is wrong. On the 24-shot bodies of our retrieval, \fsJevBGoldZero{} items (\fsJevBGoldZeroPct\% of the test set;
\fsJevBGoldZeroOfErrors\% of \jev's errors) have gold probability zero, and these items dominate its log loss. Our
pre-registered analysis found \orb's log loss lower by \fsSwapDiffLogLoss{} (CI [\fsSwapDiffLogLossLo,
\fsSwapDiffLogLossHi]); flooring \jev's probabilities at $10^{-3}$, below the rounding resolution, shrinks the
difference to \fsSwapDiffLogLossFloor. We had initially reported the log-loss difference as a calibration advantage
and withdrew it. For systems that round their outputs, only metrics that are insensitive to the rounding (ECE,
Brier) support comparisons; log loss is reported for the open models, whose release code returns unrounded
probabilities.

\subsection{Retrieval tie-breaking depends on the platform}
\label{sec:find-ties}
Our BM25 implementation ranked candidates with an unstable sort and summed term scores in hash order. Re-running
retrieval on a different processor architecture and numerical library changed the retrieved set for hundreds of
test items, all of them ties at the 24th rank. We fixed the ranking (sorted term order, stable sort with an index
tie-break) and re-scored the headline configuration post hoc: the retrieved set changes for \fsStableSetDiff{} items
(\fsStableSetDiffPct\%) and the order only for \fsStableOrderOnlyDiff{} more, but only \fsStableChanged{} predictions
change, accuracy is unchanged (\fsStableAcc\%) and the comparison with \jev{} is unaffected ($p=\fsStableVsJevP$).
The effect was negligible here, but a request manifest is the only way to know that two runs saw the same evidence.

\subsection{Selection on the test set changes significance}
\label{sec:find-selection}
Before pre-registering, we explored about forty configurations scored on the test set and found one at 93.8\% that % numbers-ok: exploratory history (claims ledger F5), not used as evidence
beat the published \jev{} result with a small $p$-value.
The repository's default configuration, by contrast, is not significantly better than \jev{} (\fsDefaultAcc\%,
$p=\fsDefaultVsJevP$), and the pre-registered configuration chosen without test data is (\fsCleanAcc\%,
$p=\fsCleanVsJevP$), under each of three test-time seeds ($p\le\fsCleanVsJevP$). The three numbers differ by less
than a point, but the significance verdict depends on how the configuration was chosen. We report only the
pre-registered result as evidence and the exploratory one as history.

\subsection{Paired and unpaired tests disagree on the headline}
\label{sec:find-paired}
The 95\% Wilson intervals of % numbers-ok: confidence level
 \orb{} and the published \jev{} result overlap ([\fsCleanWilsonLo, \fsCleanWilsonHi]
vs.\ [\fsJevPubWilsonLo, \fsJevPubWilsonHi]), which an unpaired reading would call no difference; the paired exact
test on the same items gives $p=\fsCleanVsJevP$. Comparisons of decision models on a shared test set should be
paired, which requires per-item predictions; published aggregate numbers alone do not support them.

\subsection{Small or ordered slices mislead in either direction}
\label{sec:find-slices}
The SNIPS test file is ordered by class, so its first 100 rows contain \zxSmokeSnipsClasses{} of
\zxSnipsClasses{} classes. On that slice \gj{} scores \zxSmokeSnipsNameAcc\% with label names (full test set:
\zxSmokeSnipsNameFullAcc\%) but \zxSmokeSnipsDescAcc\% with descriptions (full: \zxSmokeSnipsDescFullAcc\%). A
smoke-test slice can overstate or understate; all numbers in this paper use full test splits.

\subsection{Training-data exposure through model lineage}
\label{sec:find-exposure}
We selected PrismNLI-0.4B as \gj's NLI member because its card describes training on synthetic NLI data. The card
also states that training starts from an earlier zero-shot classifier \citep{jung-2025-prismnli-card}, whose own
training mix includes the Banking77 and MASSIVE training splits \citep{laurer-etal-2023-less}
\authorcheck{cite the pinned dataset list of the earlier classifier}. Our zero-shot numbers on those two datasets
are therefore exposed, and HWU64 is indirectly exposed because it shares utterances with MASSIVE. We replaced the
member with one whose training data are documented and exclude all seven benchmarks; the jury's internal findings
hold (\S\ref{sec:res-zeroshot}), but its absolute accuracy is lower on most datasets, including some that were never
exposed. A model card's description of the final training stage is not a statement about the data its
initialisation saw.

\subsection{Serving engines and over-length requests}
\label{sec:find-serving}
\pending{De-identified findings from deployment work, approved by the author: one model failed the parity gate on
one of \fsN{} items; a serving engine's input limit differed from the release code's and silently refused
requests the release code accepts; over-length policies differ by model and flipped one task's ranking. Each needs
its per-item evidence before it enters the text.}
